\documentclass[a4paper]{exam}
\usepackage{amsmath,amssymb,titlesec}
\usepackage[a4paper,margin=22mm]{geometry}

\titleformat{\section}{\normalfont\Large\bfseries}{}{0pt}{}
\titleformat{\subsection}{\normalfont\large\bfseries}{}{0pt}{}
\titlespacing*{\section}{1em}{1.6\baselineskip}{0.6\baselineskip}
\titlespacing*{\subsection}{2em}{1.2\baselineskip}{0.4\baselineskip}

\setlength{\parindent}{0in}
\setlength{\headheight}{18pt}

\printanswers
\title{Week 10}
\author{Radhesh Goel}
\date{\today}

\begin{document}
\maketitle

\section*{Question 1}

\subsection*{a) Temperature at the heat-flux boundary}

\begin{solution}
  At the left boundary, Fourier's law gives
  \[
    q_{\mathrm{in}} = -k\left. \frac{\partial T}{\partial x} \right|_{x=0}
  \]
  Using a first-order forward difference at time level \(m\),
  \[
    \left. \frac{\partial T}{\partial x} \right|_{x=0,t=t_m} \approx \frac{T_1^m - T_0^m}{\Delta x}
  \]
  Therefore,
  \[
    q_{\mathrm{in}} = -k\frac{T_1^m - T_0^m}{\Delta x}
  \]
  Rearranging for the boundary temperature gives
  \[
    \boxed{T_0^m=T_1^m+\frac{q_{\mathrm{in}}\Delta x}{k}}
  \]
  A positive heat flux entering the rod requires \(T_0^m > T_1^m\), consistent
  with heat flowing towards increasing \(x\)
\end{solution}

\subsection*{b) Order of accuracy}


\end{document}
